\section{Summary and Discussions}

In this work, we derive the explicit form of the density matrix of the Kitaev ground state on a torus and discuss the entanglement information encoded therein in detail. We use the spin formulation throughout this paper to address the lack of a theoretical understanding of quantum spin systems in the spin representation.

This work can be summarized in three main results. First, the ground state density matrix is composed of string operators in different edge configuration equivalent classes. The structure of the equivalence class arises from the flux-free property of the ground state, in the sense that edge configurations that can be related by loops are sorted into an equivalence class. Regardless of the fermion parity condition, an equivalence class is only hinged on the pattern of endpoints of the edge configurations therein. This structure shows the equivalence between the spin representation and the fermion representation of edge configurations. 

Second, a reduced density matrix of a simply-connected subregion is obtained. The calculation of the entanglement entropy implies that it is composed of identical blocks, and each block is the density matrix of the corresponding free-fermion system. In other words, the entanglement spectrum has an extensive degeneracy, and the distinct values in the entanglement spectrum are the ones in the entanglement spectrum of the equivalent free-fermion system. The gauge-field contribution to the entanglement entropy is actually the measure of the degeneracy in the entanglement spectrum. We also investigated the issue raised by cases with an odd number of lattice sites. Our approach solves this subtlety by borrowing a Majorana fermion from the gauge sector, and this derivation also leads to the finding of an additional 2-fold degeneracy in the entanglement entropy of the fermion sector.

Third, inspired by~\cite{PRX-ent-high}, the block-diagonal structure of the reduced density matrix is directly proven. The physical picture is based on the \textit{residual} Gauss law of the gauge theory after the environment is traced out, and it then serves as real symmetries that give good quantum numbers for block-diagonalizing the reduced density matrix. After that, the blocks are proven to be identical because the exact 1-form Wilson symmetry relates all the blocks together. Since the gauge theory of the Kitaev model is only an emergent structure in the enlarged parton space, the fractionalized fermions need to be \textit{glued} back into physical spin degrees of freedom by inserting strings that connect pairs of Majorana fermions. Lastly, a global operator defined in the sense of dimer coverings that reflects fermion parity in the subregion is found to give an additional 2-fold degeneracy when the number of lattice sites is odd, in agreement with the 2-fold degeneracy found in the Majorana fermion sector.

We would like to mention that the form of the density matrix~\eqref{equ.density-matrix} is reminiscent of the tensor network state proposed by~\cite{string-gas-TN}, in which the tensor network \textit{Ans\"atze} that approximates the Kitaev ground state is constructed by applying a tensor product operator to an initial state. The tensor product operator is composed of a \textit{loop gas} operator and a \textit{dimer gas} operator, in which the loop gas term is exactly the plaquette projector~\eqref{eq.p-proj-def}, and the dimer gas part consists of low-order terms of the summation in \eqref{equ.density-matrix}. The similarity between these two might be due to the projector nature of a pure state density matrix. If the tensor-product operator approximately captures the density matrix of the ground state, this might explain the remarkably good agreement between the results of \cite{string-gas-TN} and the exact solution.

Another remark is that the theoretical framework proposed in this work is general for every \textit{Kitaev-like} model, as long as it is exactly solvable in a way similar to the Kitaev model, in which $\mathbb{Z}_2$ fluxes have no dynamics, and the spin Hamiltonian can be completely separated into $\mathbb{Z}_2$ gauge fields and the Majorana matter parts. For example, the Kitaev model has been generalized in the way that three-spin interaction terms are included~\cite{3-spin-1, 3-spin-2, 3-spin-3, 3-spin-4, Xiang-3-spin-chiral}. The three-spin term gives rise to a time-reversal-symmetry-breaking next-nearest-neighbor Majorana hopping in the Majorana representation similar to the Haldane model~\cite{Haldane-model}, giving a chiral ground state with a non-zero Chern number. The form of the density matrix proposed in this work should also be valid for Kitaev-like models with different lattice geometries, such as the Yao-Kivelson model, which can host a chiral-spin-liquid ground state~\cite{yao-kivelson}. The lattice structure in the Yao-Kivelson model is composed of triangular and dodecagonal plaquettes, for which the corresponding plaquette operators can be defined similarly in the Kitaev model, and the pattern of the equivalence class, as well as the correspondence between strings and Majorana fermions, might still be there. If this conjecture is correct, then we could write down an explicit form of the density matrix of a chiral spin liquid formulated in terms of spin operators, which might provide a new point of view based on a spin formulation to understand the topological structure in a chiral spin liquid traditionally characterized by the Chern number in the single-particle picture.

